\documentclass[a4paper,12pt,oneside]{article}
\makeatletter
\def\input@path{{../../}}
\makeatother
\usepackage{discrete-practice}

% Источник условий: https://nebulous-name-0d3.notion.site/2-19-03-2024-23-59-30d899a19f54401cbc9602193d10568d?pvs=25
\begin{document}

\HomeworkHeader{2}{2}{Графы. Карта Европы}

Граф\footnote{Простой неориентированный граф.} Европы $G^*=\langle V,E\rangle$
определён следующим образом. Каждая вершина $v\in V$ соответствует европейской
стране\footnote{Можно использовать \href{http://simple.wikipedia.org/wiki/List_of_European_countries}{список европейских стран}.}.
Две вершины смежны ($\{u,v\}\in E$), если соответствующие им страны имеют
сухопутную границу. Пусть $G$ будет максимальной компонентой связности $G^*$.

\begin{numlist}
  \item Постройте $G^*$ с минимальным пересечением рёбер.
  \item Вычислите
    \[
      |V|,\; |E|,\; \delta(G)=\min\deg(v),\; \Delta(G)=\max\deg(v),\;
      \mathrm{rad}(G),\; \mathrm{diam}(G),\; \mathrm{center}(G)
    \]
    и цикломатическое число.
  \item Вычислите хроматическое число $\chi(G)$.
  \item Найдите максимальную клику $Q\subseteq V$ графа $G$.
  \item Найдите простой цикл $W$, который содержит все вершины графа $G$.
  \item Найдите цикл $U$, который содержит все рёбра графа $G$.
  \item Перечислите все компоненты вершинной двусвязности и постройте граф
    блоков и точек сочленения графа $G^*$.
  \item Перечислите все компоненты рёберной двусвязности и постройте граф
    компонент рёберной двусвязности графа $G^*$.
  \item Добавьте веса на рёбра, равные расстоянию\footnote{Можно выбрать
    геодезическое или автомобильное расстояние.} между столицами двух стран.
    Найдите минимальное остовное дерево $T$ для $G$.
  \item Закодируйте граф $T$ кодом Прюфера и бинарным кодом.
\end{numlist}

\end{document}
